\section{Compounding: What Really Grows When a Company Grows?}

\subsection{Beyond Revenue: The Accumulation of Assets}

The attraction of a successful business is not merely that it earns money today. It is that the work done today can make tomorrow's work more productive.

Consider a specialist engineering firm in its 1st year. The team spends weeks learning how to solve a particular class of difficult client problems. It documents its approach, builds diagnostic tools, learns which errors recur, and develops a reputation. In the 2nd year, the same kind of project takes less time and produces fewer mistakes. Satisfied clients introduce the firm to peers. The firm hires an engineer who learns through the accumulated methods rather than rediscovering everything. Soon the firm can serve more clients with greater consistency.

The 1st year's work has produced more than revenue. It has produced assets: knowledge, tools, relationships, evidence, and a repeatable operating method.

This is a business form of compounding.

\subsection{Organizational Feedback Loops and Owner Earnings}

Financial compounding reinvests a return so that subsequent returns accrue on a larger base. Organizational compounding is more complicated. Sometimes experience reduces costs. Sometimes trust improves conversion. Sometimes product data improves quality. Sometimes a better brand attracts stronger employees, who create better products, which reinforce the brand. Sometimes these loops reverse: rapid growth damages quality, weak quality harms reputation, and recruitment becomes harder.

A founder should ask what is accumulating after each customer, employee, project, or transaction. If nothing accumulates but cash, the business may still be attractive. But if every job must be sold and delivered almost from scratch, growth may remain tied to additional effort. If each project makes subsequent projects easier to win and execute, the company is acquiring a capability that can outlast a single transaction.

Warren Buffett has repeatedly drawn attention to businesses that can earn strong returns while needing relatively little additional capital to maintain or expand their operations. In his 1986 Berkshire Hathaway shareholder letter, he distinguished reported accounting earnings from the resources owners could economically extract without weakening a business.\footnote{Warren Buffett, Berkshire Hathaway Shareholder Letter (1986).} The underlying insight applies to founders: profit that must be continuously reinvested merely to preserve the current business is different from surplus that can be invested in genuinely new opportunities.

The quality of compounding depends on 3 things: the rate of return on incremental resources, the amount that can be usefully reinvested, and the duration for which attractive opportunities persist. A business that earns high returns on a tiny amount of capital but cannot expand may be a wonderful cash-generating enterprise without becoming a giant. A business that can reinvest large sums at attractive returns for decades has a much larger potential economic footprint.

Compounding does not happen simply because a spreadsheet assumes 30\% growth. Growth must be earned by solving additional customer problems, reaching more buyers, improving capabilities, or investing successfully. Competition and diminishing returns usually make that harder over time.

\textbf{This is why the highest aspiration of a founder should not be perpetual hypergrowth as a slogan. It should be to build an engine that converts today's efforts and resources into tomorrow's productive capabilities.}
